\documentclass[11pt]{article}
\usepackage[margin=1in]{geometry}
\usepackage{amsmath,amssymb,amsthm}
\usepackage{hyperref}
\usepackage{enumitem}

\newtheorem{theorem}{Theorem}
\newtheorem{lemma}[theorem]{Lemma}
\newtheorem{corollary}[theorem]{Corollary}
\theoremstyle{definition}
\newtheorem{definition}[theorem]{Definition}
\theoremstyle{remark}
\newtheorem{remark}[theorem]{Remark}

\usepackage{xurl}
\let\oldus\_
\renewcommand{\_}{\oldus\allowbreak}

\title{A proof of Conjecture 00000007000\\[2pt]
\large A permutation pair that preserves the box norm but changes the scattering spectrum}
\date{}

\begin{document}
\sloppy
\maketitle

\section{The conjecture}

Conjecture 00000007000 reads:

\begin{quote}
\emph{Definition:} Box norms and scattering spectra are two layers. \emph{Conjecture:} There exist
two functions with identical box norms but different scattering spectra, and the separation is
realized by an explicit permutation pair with the same norms but different phase.
\end{quote}

It is an existential statement. We prove it with two explicit real functions on the grid
$\mathbb{Z}/4\times\mathbb{Z}/4$.

\section{Reading and definitions}

\begin{definition}[Box norm]
For finite sets $X,Y$ and $f:X\times Y\to\mathbb{R}$, the (Gowers) box norm is
\[
  \|f\|_{\square}=\Bigl(\mathbb{E}_{x,x'\in X}\,\mathbb{E}_{y,y'\in Y}\,
  f(x,y)f(x,y')f(x',y)f(x',y')\Bigr)^{1/4}.
\]
The inner quantity equals $|X|^{-2}|Y|^{-2}\sum_{x,x'}\bigl(\sum_y f(x,y)f(x',y)\bigr)^2\ge0$.
For a function of two variables this is the box norm; we prove its invariance for every
permutation pair.
\end{definition}

\begin{definition}[Permutation pair]
A permutation pair is a pair $(\sigma,\tau)$ of a permutation $\sigma$ of $X$ and a permutation
$\tau$ of $Y$. It acts by $f\mapsto f\circ(\sigma\times\tau)$, i.e.\ $(x,y)\mapsto
f(\sigma x,\tau y)$.
\end{definition}

\begin{definition}[Scattering spectrum, phase]
On $G=\mathbb{Z}/4\times\mathbb{Z}/4$ the Fourier transform is
$\hat f(\xi)=\sum_{(x,y)\in G}f(x,y)\,e^{-2\pi i(\xi_1x+\xi_2y)/4}$. The scattering spectrum is the
diffraction intensity (structure factor) $I_f(\xi)=|\hat f(\xi)|^2$, and the phase at $\xi$ is
$\hat f(\xi)/|\hat f(\xi)|$.
\end{definition}

In Lean the character is written $i^{k}$ for $k\in\mathbb{Z}/4$ (using the representative
$k\in\{0,1,2,3\}$); the lemma \texttt{exp\_eq} proves $i^{k}=e^{2\pi ik/4}$.

\section{Results}

\begin{lemma}[Permutation pairs preserve the box norm]\label{lem:inv}
For all $f$ and every permutation pair, $\|f\circ(\sigma\times\tau)\|_\square=\|f\|_\square$.
\end{lemma}

\begin{proof}
The sum over $(x,x',y,y')$ is reindexed by the bijection
$(x,x',y,y')\mapsto(\sigma x,\sigma x',\tau y,\tau y')$ of $X\times X\times Y\times Y$.
\end{proof}

\paragraph{The explicit pair.} Let $f(x,y)=1$ if $x\in\{0,1\}$ and $y=0$, and $f=0$ otherwise.
Let $\sigma=(1\ 2)$ and $\tau=(0\ 1)$ (transpositions of $\mathbb{Z}/4$), and $g=f\circ(\sigma\times
\tau)$. Then $g(x,y)=1$ iff $x\in\{0,2\}$ and $y=1$.

\begin{theorem}[Conjecture 00000007000]\label{thm:main}
\begin{enumerate}[nosep]
\item $\sigma\ne\mathrm{id}$, $\tau\ne\mathrm{id}$, and $g=f\circ(\sigma\times\tau)$;
\item identical box norms: $\|g\|_\square=\|f\|_\square>0$;
\item different scattering spectra: $I_f(1,0)=2$ while $I_g(1,0)=0$;
\item different phase: $I_f(0,1)=I_g(0,1)=4$, but the phase of $\hat f(0,1)$ is $1$ and that of
$\hat g(0,1)$ is $-i$.
\end{enumerate}
\end{theorem}

\begin{proof}
(2) The equality is Lemma~\ref{lem:inv}. The box-norm sum of $f$ is
$\sum_{x,x'}(\sum_yf(x,y)f(x',y))^2=4$, so $\|f\|_\square=(4/256)^{1/4}>0$.

(3) $\hat f(1,0)=1+e^{-2\pi i/4}=1-i$, so $I_f(1,0)=2$. $\hat g(1,0)=1+e^{-2\pi i\cdot 2/4}=1-1=0$.

(4) $\hat f(0,1)=2$ (both points have $y=0$), while $\hat g(0,1)=2e^{-2\pi i/4}=-2i$. Both have
modulus $2$; the phases are $1$ and $-i$.
\end{proof}

\begin{remark}
A translation of $\mathbb{Z}/4$ is a permutation that changes only the phases of $\hat f$, never the
intensities. The permutation $\sigma=(1\ 2)$ is not affine on $\mathbb{Z}/4$, which is why it can
change the intensity at $(1,0)$. For the displayed function $f$, which is supported on the single column
$y=0$, the transposition $\tau=(0\ 1)$ only moves that column to $y=1$ and so multiplies $\hat f$ by a phase in
the $y$ direction; $\tau$ is not a translation, and for general functions it does change intensities. (This remark is
explanatory only; it is not part of the Lean formalization.)
\end{remark}

\section{Lean formalization}

The file \texttt{lean/Conjecture7000/Basic.lean} (Lean 4, Mathlib v4.33.1, namespace
\texttt{Submission00000007000}) defines \texttt{boxSum}, \texttt{boxNorm} (with real power $1/4$),
\texttt{permute}, \texttt{e4}, \texttt{fourier}, \texttt{intensity} (\texttt{Complex.normSq} of the
Fourier coefficient) and \texttt{phase}. The main statements are \texttt{boxSum\_eq\_sq},
\texttt{boxSum\_permute}, \texttt{boxNorm\_permute} (Lemma~\ref{lem:inv}), \texttt{exp\_eq},
\texttt{g\_eq}, \texttt{fourier\_f\_10}, \texttt{fourier\_g\_10}, \texttt{fourier\_f\_01},
\texttt{fourier\_g\_01}, \texttt{phase\_f\_01}, \texttt{phase\_g\_01}, \texttt{boxSum\_f}, and
\begin{verbatim}
theorem conjecture7000 :
    exists (f g : G -> R) (s t : Equiv.Perm (ZMod 4)),
      s != 1 /\ t != 1 /\ g = permute s t f /\
      boxNorm g = boxNorm f /\ 0 < boxNorm f /\
      intensity f != intensity g /\ intensity f (1, 0) = 2 /\
      intensity g (1, 0) = 0 /\
      intensity f (0, 1) = intensity g (0, 1) /\
      phase f (0, 1) != phase g (0, 1)
\end{verbatim}
(written here in ASCII; the Lean source uses Unicode symbols and the names $\sigma,\tau$).

\paragraph{Axiom audit.} \texttt{lake build} succeeds. \texttt{\#print axioms} for
\texttt{conjecture7000}, \texttt{boxNorm\_permute} and \texttt{exp\_eq} reports only
\texttt{propext}, \texttt{Classical.choice} and \texttt{Quot.sound}. There is no \texttt{sorry}, no
\texttt{native\_decide} and no added axiom.

\section{Scope}

``Box norm'' is read as the Gowers box norm of a real function of two variables; ``scattering
spectrum'' as the diffraction intensity $|\hat f|^2$ on the finite grid; ``permutation pair'' as a
pair of coordinate permutations; ``phase'' as the unit phase of a Fourier coefficient. Under this
reading both the intensity and the phase change. We do not treat the wavelet scattering transform
(Mallat) or quantum-mechanical scattering matrices.

\end{document}
